\usepackage[T1]{fontenc} %pacote inicial fonte
\usepackage[margin=1in]{geometry} % para margem principalmente, mas tambem formatar texto
\usepackage{xcolor} % para cor de texto
\usepackage{lipsum} % para texto aleatorio
\usepackage[latin,brazilian]{babel} % para linguas entendidas pelo texto
\usepackage{graphicx} % para imagens
\usepackage{array} %para tabelas
\usepackage{booktabs} %para tabelas tambem, so que mais completo
\usepackage{amsmath} % para matematica
\usepackage{mathcomp}
\usepackage{amsfonts} % para formatacao de texto matematico
\usepackage{amssymb}
\usepackage{url}
\usepackage{hyperref}
\usepackage{setspace}
\usepackage[parfill]{parskip} % para espaçamento no texto entre paragrafos
\usepackage{natbib} % para citacoes tambem exige que os arquivos bib sejam excluidos
%\usepackage[style=authoryear]{biblatex} % tomar cuidado que biblatex não tem suporte para o latin e pode exigir recompilacao do bib
\usepackage{csquotes}
%\addbibresource{learnlatex.bib} %file of reference
